\section{Related Literature}
\label{sec:literature}

Our paper addresses four gaps in the literature on bid rigging and
procurement design.

\paragraph{Gap 1: Coordination tests require \emph{ad hoc} bidder
groupings.}
\citet{bajari2003deciding} develop powerful tests---exchangeability of
bid functions and conditional independence of bid residuals---but
require the researcher to partition bidders into ``suspected
colluders'' and ``competitive fringe'' before estimation.
\citet{porter1993detection} and \citet{porter1999ohio} detect bid
rotation in Ohio school milk markets; \citet{conley2016detecting}
detect bidder groups from co-bidding patterns without bid values;
\citet{kawai2019using} use the margin of victory to rule out
cost-differential explanations;
\citet{clark2021collusion} study collusion dynamics in long-lived
cartels; \citet{schurter2020identification} addresses identification
in first-price auctions with collusion.
In each case, the partition is constructed \emph{ex post} or by
assumption. Our FL classification provides a data-driven,
\emph{ex ante} partition into suspected cover bidders and genuine
competitors, grounded in a participation anomaly that is difficult
to rationalize as competitive behavior.
Notably, \citet{porter1993detection} show that bid-rigging cartels
strategically self-select into competitive markets---a pattern we
document in our network-split analysis
(Section~\ref{sec:results_network}) and event study
(Section~\ref{sec:additional_id}).

\paragraph{Gap 2: Existing screens require bid-level data.}
\citet{imhof2019detecting} and \citet{imhof2018screening} develop
machine-learning screens using tender-level features (bid variance,
coefficient of variation, kurtosis) that achieve high accuracy against
known cartels. \citet{abrantesmetz2006a} propose variance screens;
\citet{harrington2008detecting} provides a comprehensive survey of
detection methods; \citet{hüschelrath2014cartel} review enforcement
complementarities between screening and leniency;
\citet{chassang2022robust} design robust screens for non-competitive
bidding; \citet{wallimann2023machine} extend ML detection to incomplete
cartels. The strategic interaction between cartel adaptation and
enforcement implies that detection methods must evolve continuously,
as cartels adjust behavior to evade known screens
\citep{harrington2008detecting}.
These approaches all operate at the \emph{tender} level and
require granular bid values---data that are unavailable or unreliable
in many procurement systems. Our FL screen operates at the \emph{firm}
level using only participation and outcome records, enabling a
two-stage workflow: FL screening first, then bid-level analysis of
flagged tenders.%
\footnote{Our Imhof comparison uses a CV median split as a coarser
approximation of the full \citet{imhof2019detecting} implementation;
see Section~\ref{sec:results_detection} for caveats on this
comparison.}

\paragraph{Gap 3: Cover bidding is documented \emph{ex post}, not
detected \emph{ex ante}.}
\citet{porter1999ohio} document complementary bidding;
\citet{pesendorfer2000study} analyzes coordinated bids in school milk
auctions; \citet{asker2010leniency} describes internal role allocation
in a stamp cartel; \citet{marshall2012economics} provide a
comprehensive treatment of cover-bidding organization;
\citet{athey2011comparing} compare auction formats under collusion,
showing that bid distributions differ systematically between
competitive and collusive regimes. In all cases,
cover bidders are identified from prosecuted cases---after the cartel
has been discovered by other means. We identify cover bidders
systematically from participation data, turning a \emph{post hoc}
observation into a \emph{proactive} detection tool.

\paragraph{Gap 4: Limited evidence from developing economies.}
The bid-rigging literature focuses predominantly on developed
economies, yet enforcement is most challenging where competition
authorities have limited resources and procurement data systems are
nascent \citep{sanchezgraells2019screening}. Brazil's BEC
platform---4.5 million tender-items, 40 million bids, 41,000 firms
over 11 years, with both sealed-bid (convite) and electronic auction
(preg\~ao) modalities---enables analyses infeasible in smaller
datasets and provides a test bed for screens designed to work with
minimal data infrastructure.
